\documentclass[UTF8,a4paper,11pt]{ctexart}
\usepackage[margin=22mm,headheight=14pt,footskip=12mm]{geometry}
\usepackage{amsmath,amssymb,array,longtable,booktabs,tikz,needspace,enumitem,fancyhdr,hyperref}
\usepackage{microtype}
\hypersetup{colorlinks=false,pdfborder={0 0 0},pdfauthor={Good Learning},bookmarks=false}
\newfontfamily\circledfont{Segoe UI Symbol}
\newcolumntype{P}[1]{>{\raggedright\arraybackslash}p{#1}}
\setlength{\LTleft}{0pt}\setlength{\LTright}{0pt}
\setlength{\parindent}{0pt}\setlength{\parskip}{.65em}
\linespread{1.18}
\raggedbottom
\setcounter{secnumdepth}{0}
\setlist{leftmargin=2.2em,itemsep=.4em,topsep=.5em}
\ctexset{section={format=\Large\bfseries\sffamily,beforeskip=1.35em,afterskip=.65em},
subsection={format=\large\bfseries\sffamily,beforeskip=1em,afterskip=.45em}}
\newcommand{\chaptertitle}[1]{\Needspace{8\baselineskip}\noindent{\LARGE\bfseries\sffamily #1}\par\vspace{4mm}\hrule\vspace{6mm}}
\newcommand{\CompactQuestion}[2]{\par\noindent\textbf{#1.}\ #2\par\vspace{.8em}}
\newcommand{\PracticeQuestion}[3]{\par\Needspace{4\baselineskip}\noindent\begin{minipage}{\linewidth}\textbf{#1.}\ #2\par\noindent\rule{0pt}{#3mm}\par\end{minipage}\par}
\newcommand{\CompactQuestionColon}[2]{\par\noindent\textbf{#1：}\ #2\par\vspace{.8em}}
\newcommand{\PracticeQuestionColon}[3]{\par\Needspace{4\baselineskip}\noindent\begin{minipage}{\linewidth}\textbf{#1：}\ #2\par\noindent\rule{0pt}{#3mm}\par\end{minipage}\par}
\pagestyle{fancy}\fancyhf{}\fancyfoot[C]{\thepage}\renewcommand{\headrulewidth}{0pt}
\begin{document}
\thispagestyle{empty}
\noindent{\sffamily\footnotesize 每日课程}\hfill{\sffamily\footnotesize 学习导图}\par
\vspace{7mm}
\noindent{\sffamily\fontsize{25}{30}\selectfont\bfseries 第1章 预备知识（宽松排版型）}\par
\vspace{4mm}\noindent\rule{\linewidth}{.85pt}\par\vspace{6mm}
\begin{center}
\begin{tikzpicture}[x=1cm,y=1cm,>=stealth,every node/.style={font=\sffamily}]
\path[use as bounding box] (-8.1,-18.7) rectangle (8.1,.4);
\node[font=\sffamily\Large\bfseries] at (0,-.6) {第1章 预备知识};
\node[font=\sffamily\small,text=black!55] at (0,-1.25) {从分布函数到条件期望};
\node[anchor=north west,font=\sffamily\fontsize{22}{24}\selectfont\bfseries,text=black!50] at (-7.35,-2.35) {01};
\node[anchor=north west,text width=5.8cm,align=left,font=\sffamily\large\bfseries] at (-6.199999999999999,-2.3200000000000003) {概率空间与分布函数};
\node[anchor=north west,text width=5.9cm,align=left,font=\sffamily\small] at (-6.199999999999999,-3.0700000000000003) {1.1—1.2\\[.25em]分布函数三性质\\[.25em]单点概率与区间概率};
\draw[black!35,line width=.4pt] (-7.35,-5.050000000000001) -- (-0.2999999999999998,-5.050000000000001);
\node[anchor=north west,font=\sffamily\fontsize{22}{24}\selectfont\bfseries,text=black!50] at (0.85,-2.35) {02};
\node[anchor=north west,text width=5.8cm,align=left,font=\sffamily\large\bfseries] at (2.0,-2.3200000000000003) {数字特征与矩母函数};
\node[anchor=north west,text width=5.9cm,align=left,font=\sffamily\small] at (2.0,-3.0700000000000003) {1.3 前半\\[.25em]ψ(t) 的导数给出矩\\[.25em]独立和变乘积};
\draw[black!35,line width=.4pt] (0.85,-5.050000000000001) -- (7.8999999999999995,-5.050000000000001);
\node[anchor=north west,font=\sffamily\fontsize{22}{24}\selectfont\bfseries,text=black!50] at (0.85,-7.7) {03};
\node[anchor=north west,text width=5.8cm,align=left,font=\sffamily\large\bfseries] at (2.0,-7.67) {特征函数};
\node[anchor=north west,text width=5.9cm,align=left,font=\sffamily\small] at (2.0,-8.42) {1.3 后半\\[.25em]五个必背公式\\[.25em]唯一性定理定分布};
\draw[black!35,line width=.4pt] (0.85,-10.4) -- (7.8999999999999995,-10.4);
\node[anchor=north west,font=\sffamily\fontsize{22}{24}\selectfont\bfseries,text=black!50] at (-7.35,-7.7) {04};
\node[anchor=north west,text width=5.8cm,align=left,font=\sffamily\large\bfseries] at (-6.199999999999999,-7.67) {收敛性与条件期望};
\node[anchor=north west,text width=5.9cm,align=left,font=\sffamily\small] at (-6.199999999999999,-8.42) {1.4—1.5\\[.25em]全期望与条件方差\\[.25em]Wald 公式的条件};
\draw[black!35,line width=.4pt] (-7.35,-10.4) -- (-0.2999999999999998,-10.4);
\node[anchor=north west,font=\sffamily\fontsize{22}{24}\selectfont\bfseries,text=black!50] at (-7.35,-13.05) {05};
\node[anchor=north west,text width=5.8cm,align=left,font=\sffamily\large\bfseries] at (-6.199999999999999,-13.020000000000001) {独立练习};
\node[anchor=north west,text width=5.9cm,align=left,font=\sffamily\small] at (-6.199999999999999,-13.770000000000001) {A1—A8\\[.25em]限时 60 分钟\\[.25em]先做后对答案};
\draw[black!35,line width=.4pt] (-7.35,-15.75) -- (-0.2999999999999998,-15.75);
\node[anchor=north west,font=\sffamily\fontsize{22}{24}\selectfont\bfseries,text=black!50] at (0.85,-13.05) {06};
\node[anchor=north west,text width=5.8cm,align=left,font=\sffamily\large\bfseries] at (2.0,-13.020000000000001) {判分与反馈};
\node[anchor=north west,text width=5.9cm,align=left,font=\sffamily\small] at (2.0,-13.770000000000001) {标注不确定题\\[.25em]拍照发我\\[.25em]据此调整第 2 章密度};
\draw[black!35,line width=.4pt] (0.85,-15.75) -- (7.8999999999999995,-15.75);
\draw[->,black!55,line width=.8pt] (-.12,-3.7) -- (.55,-3.7);
\draw[->,black!55,line width=.8pt] (7.85,-5.45) -- (7.85,-7.2);
\draw[->,black!55,line width=.8pt] (.52,-9.05) -- (-.15,-9.05);
\draw[->,black!55,line width=.8pt] (-7.85,-10.8) -- (-7.85,-12.55);
\draw[->,black!55,line width=.8pt] (-.12,-14.4) -- (.55,-14.4);
\node[font=\sffamily\small,text=black!55] at (0,-18.1) {教材 1.1—1.5 · 笔记第 1—2 页};
\end{tikzpicture}
\end{center}
\clearpage

\chaptertitle{第1章 预备知识（D001）}
教材：张波、商豪、邓军《应用随机过程》（第6版），中国人民大学出版社，2023。对应 1.1—1.5（第 001—031 页）。本章不是“复习一下就行”，第 1.3 节的矩母函数、特征函数与第 1.5 节的条件期望是后面泊松过程、马尔可夫链、鞅的工具；2023 级期末第 1 题就只考特征函数。

\textbf{今天的时间分配（2.5 小时）}：概念与例题 80 分钟，填步题与变式 60 分钟，独立练习 40 分钟，订正与回顾 10 分钟。

\Needspace{6\baselineskip}
\section{一、开课前 10 分钟：自检、节奏与用法}
\textbf{前置自检（3 题，限 3 分钟，不确定就标记）}

\begin{enumerate}
\item 若 $X$ 取 $0,1$ 各 $\frac12$，$E(X)$ 与 $\mathrm{Var}(X)$ 分别是多少？
\item 若 $X,Y$ 相互独立，$E(XY)$ 等于什么？
\item 正态分布 $N(\mu,\sigma^2)$ 的密度函数写得出吗？
\end{enumerate}
三题都会，就按下面的节奏直接开始；有一题卡住，先补对应的高数或概率课本节，再回到单元 1。本章表面是“复习”，实际新增了两件工具：矩母函数与特征函数（1.3 节）、条件期望（1.5 节）。第 3 章泊松过程、第 5 章马尔可夫链、第 6 章鞅全部依赖它们；2023 级期末卷第一道大题就考特征函数，属于必须拿满分的题。

\textbf{今晚节奏（2.5 小时）}

{\small\renewcommand{\arraystretch}{1.38}\setlength{\tabcolsep}{3.5pt}
\begin{longtable}{P{0.17\linewidth}P{0.48\linewidth}P{0.28\linewidth}}
\toprule
\bfseries 时间 & \bfseries 做什么 & \bfseries 产出 \\
\midrule
\endfirsthead
\toprule
\bfseries 时间 & \bfseries 做什么 & \bfseries 产出 \\
\midrule
\endhead
0—10 分钟 & 自检 3 题，扫一遍四个单元标题 & 知道哪一节要放慢 \\
\addlinespace[.35em]
10—90 分钟 & 单元 1—4：读概念、跟一遍完整例题 & 笔记：五个特征函数、条件方差公式 \\
\addlinespace[.35em]
90—150 分钟 & 填步题 R1—R4、变式 V1—V4 & 每题写出完整步骤 \\
\addlinespace[.35em]
150—190 分钟 & 独立练习 A1—A8（限时） & 答在宽松版的空白处 \\
\addlinespace[.35em]
190—210 分钟 & 对答案，标出不确定的题 & 反馈清单（拍照发我） \\
\addlinespace[.35em]
\bottomrule
\end{longtable}
}
\textbf{两版 PDF 怎么用}：紧凑版用来读和复习，宽松版题目下留了手写空白，用来作答；先做完独立练习再看第六部分答案，不要边做边对。答案只给关键步骤，写不出时回到对应单元的“完整例题”看方法，不要直接抄答案。

\clearpage
\Needspace{6\baselineskip}
\section{单元 1 概率空间与随机变量（1.1—1.2）}
\textbf{概念与条件}

\begin{itemize}
\item 概率空间记作 $(\Omega,\mathcal{F},P)$：$\Omega$ 是样本空间，$\mathcal{F}$ 是对补、可列并封闭的事件域，$P$ 满足非负、$P(\Omega)=1$、可列可加。
\item 随机变量 $X$ 是 $\Omega$ 上的实值函数，且对任意实数 $x$ 有 $\{\omega:X(\omega)\le x\}\in\mathcal{F}$。这个“可测”条件保证概率 $P(X\le x)$ 有意义，计算时几乎不用验证，但要知道它为什么在。
\item 分布函数 $F(x)=P(X\le x)$ 的三条性质：单调不减、右连续、$F(-\infty)=0$、$F(+\infty)=1$。反过来，满足这三条的函数一定是某个随机变量的分布函数。
\item 由分布函数求区间概率：$P(a<X\le b)=F(b)-F(a)$。注意左端是开区间，$P(X=a)=F(a)-F(a-0)$，连续型随机变量单点概率为 0。
\item 期望用 Riemann–Stieltjes 积分统一写法：$E(X)=\int_{-\infty}^{+\infty}x\,dF(x)$。离散型退化为求和，连续型退化为 $\int x f(x)dx$。
\end{itemize}
\textbf{完整例题（教材例 1.2.2 的处理方式）}

设 $X$ 等可能地取 $0$ 与 $1$，求 $X$ 的分布函数。

第一步，写出分布列：$P(X=0)=P(X=1)=\frac12$。 第二步，按 $x$ 落在哪个区间分段讨论，因为分布函数是“累积到 $x$”的概率： 当 $x<0$ 时没有质量，$F(x)=0$； 当 $0\le x<1$ 时只有 $X=0$ 这一点，$F(x)=P(X=0)=\frac12$； 当 $x\ge 1$ 时两点都包含，$F(x)=1$。 第三步，写成分段函数并检查：在 $x=0,1$ 处函数值等于右极限，所以右连续。 结论：$F(x)=0\ (x<0)$，$F(x)=\frac12\ (0\le x<1)$，$F(x)=1\ (x\ge 1)$。

\Needspace{4\baselineskip}
\subsection{举一反三}
\PracticeQuestion{R1}{设 $X$ 等可能地取 $-1,0,1$，求 $F(x)$，并求 $P(0<X\le 1)$。}{42}
\PracticeQuestion{V1}{设 $X$ 的分布列为 $P(X=0)=\frac14$、$P(X=1)=\frac12$、$P(X=2)=\frac14$，求 $F(x)$ 与 $P(0<X\le 2)$；再问：$P(X=1)$ 能否用 $F$ 表示？}{30}
\PracticeQuestion{M1}{判断下列函数是否为分布函数，并说明理由：$F(x)=0\ (x<0)$，$F(x)=x\ (0\le x<1)$，$F(x)=1\ (x\ge 1)$ 与 $G(x)=0\ (x\le 0)$，$G(x)=\frac12\ (x>0)$。}{40}
\clearpage
\Needspace{6\baselineskip}
\section{单元 2 数字特征与矩母函数（1.3 前半）}
\textbf{概念与条件}

\begin{itemize}
\item $E(X)=\int x\,dF(x)$；$\mathrm{Var}(X)=E(X^2)-[E(X)]^2$；$\mathrm{Cov}(X,Y)=E(XY)-E(X)E(Y)$。
\item 矩母函数 $\psi(t)=E(e^{tX})=\int e^{tx}dF(x)$，要求在 $t=0$ 的某个邻域内有限。它存在的条件比特征函数强，指数分布、正态分布有，但重尾分布可能没有。
\item 关键性质：$\psi(0)=1$；若可在 $t=0$ 处求导，则 $\psi^{(k)}(0)=E(X^k)$，于是 $E(X)=\psi'(0)$、$E(X^2)=\psi''(0)$、$\mathrm{Var}(X)=\psi''(0)-[\psi'(0)]^2$。
\item 独立时 $\psi_{X+Y}(t)=\psi_X(t)\psi_Y(t)$。这条是把“和的分布”变成“乘积”的桥梁。
\item 注意：矩母函数存在时与分布一一对应（唯一性），这是与特征函数相同的性质。
\end{itemize}
\textbf{完整例题}

设 $X\sim B(n,p)$，求 $\psi_X(t)$，并用它求 $E(X)$ 与 $\mathrm{Var}(X)$。

第一步，按定义展开：$\psi(t)=E(e^{tX})=\sum_{k=0}^{n}e^{tk}\binom{n}{k}p^k(1-p)^{n-k}$。 第二步，合并同底数：$e^{tk}p^k=(pe^{t})^k$，于是 $\psi(t)=\sum_{k}\binom{n}{k}(pe^{t})^k(1-p)^{n-k}=(pe^{t}+1-p)^n$，用的是二项式定理。 第三步，求导：$\psi'(t)=n(pe^{t}+1-p)^{n-1}pe^{t}$，代入 $t=0$ 得 $E(X)=np$。 第四步，再求导：$\psi''(t)=n(n-1)(pe^{t}+1-p)^{n-2}(pe^{t})^2+n(pe^{t}+1-p)^{n-1}pe^{t}$，代入 $t=0$ 得 $E(X^2)=n(n-1)p^2+np$，所以 $\mathrm{Var}(X)=np(1-p)$。 每一步都用到了“在 $t=0$ 处取值”这个动作，这是矩母函数题固定套路。

\Needspace{4\baselineskip}
\subsection{举一反三}
\PracticeQuestion{R2}{设 $X\sim \mathrm{Exp}(\lambda)$，其密度为 $\lambda e^{-\lambda x}\ (x>0)$。先算 $\psi(t)=\frac{\lambda}{\lambda-t}\ (t<\lambda)$，再求 $E(X)$ 与 $\mathrm{Var}(X)$。}{38}
\PracticeQuestion{V2}{设 $X\sim P(\lambda)$（泊松分布），求 $\psi_X(t)$，并求 $E(X)$ 与 $\mathrm{Var}(X)$；与指数分布的结果比较，说明为什么泊松分布的 $\psi$ 在任意 $t$ 都有定义。}{42}
\PracticeQuestion{M2}{设 $X,Y$ 独立，$X\sim B(n,p)$、$Y\sim B(m,p)$。用矩母函数说明 $X+Y\sim B(n+m,p)$，并指出这里用到了哪条性质。}{34}
\clearpage
\Needspace{6\baselineskip}
\section{单元 3 特征函数（1.3 后半）}
\textbf{概念与条件}

\begin{itemize}
\item 特征函数 $\varphi(t)=E(e^{itX})=\int e^{itx}dF(x)$，对任何分布都存在，且 $\varphi(0)=1$、$|\varphi(t)|\le 1$、在 $\mathbb{R}$ 上一致连续。
\item 反号关系：$\varphi_{-X}(t)=\overline{\varphi_X(t)}$；线性变换：$\varphi_{aX+b}(t)=e^{ibt}\varphi_X(at)$。
\item 独立和：$\varphi_{X+Y}(t)=\varphi_X(t)\varphi_Y(t)$。
\item 唯一性定理：特征函数与分布函数相互唯一决定。判断“和服从什么分布”时，算出乘积后认出一个已知特征函数即可，不必求反演积分。
\item 矩公式：$\varphi^{(k)}(0)=i^kE(X^k)$，即 $E(X)=\frac{\varphi'(0)}{i}$、$E(X^2)=-\varphi''(0)$。这条在求高阶矩时比求导矩母函数更稳。
\end{itemize}
\textbf{常见特征函数（必须背下来）}：两点分布 $pe^{it}+1-p$；二项分布 $(pe^{it}+1-p)^n$；泊松分布 $e^{\lambda(e^{it}-1)}$；指数分布 $\left(1-\frac{it}{\lambda}\right)^{-1}$；正态分布 $e^{i\mu t-\frac12\sigma^2t^2}$。

\textbf{完整例题（2023 级期末卷第一道大题）}

设 $X\sim B(n,p)$，$Y\sim B(m,p)$，且 $X$ 与 $Y$ 相互独立。求 $X$ 的特征函数与 $X+Y$ 的分布。

第一步，写出 $X$ 的特征函数：$\varphi_X(t)=\sum_k e^{itk}\binom{n}{k}p^k(1-p)^{n-k}=(pe^{it}+1-p)^n$。 第二步，同理 $\varphi_Y(t)=(pe^{it}+1-p)^m$。 第三步，用独立性把和的特征函数拆成乘积：$\varphi_{X+Y}(t)=\varphi_X(t)\varphi_Y(t)=(pe^{it}+1-p)^{n+m}$。 第四步，认出这是参数为 $n+m$、$p$ 的二项分布的特征函数，由唯一性定理得 $X+Y\sim B(n+m,p)$。 要点：第三步的“独立”不能省；第四步必须点明唯一性定理，否则只算出一个函数，没有回答“分布是什么”。

\Needspace{4\baselineskip}
\subsection{举一反三}
\PracticeQuestion{R3}{设 $X\sim P(\lambda_1)$、$Y\sim P(\lambda_2)$ 且独立，求 $\varphi_{X+Y}(t)$ 并写出 $X+Y$ 的分布。}{42}
\PracticeQuestion{V3}{设 $X\sim N(0,1)$，用 $\varphi_X(t)=e^{-t^2/2}$ 求 $E(X^2)$ 与 $E(X^4)$，并说明为什么不能直接对 $\varphi$ 求一阶导数得到 $E(X)$。}{34}
\PracticeQuestion{M3}{判断并说明理由：若两个随机变量的特征函数相同，则它们的分布函数相同；若两个随机变量的矩母函数相同，是否也一定相同？}{34}
\clearpage
\Needspace{6\baselineskip}
\section{单元 4 收敛性与条件期望（1.4—1.5）}
\textbf{概念与条件}

\begin{itemize}
\item 四种收敛：依概率 $X_n\xrightarrow{P}X$（对任意 $\varepsilon>0$，$P(|X_n-X|>\varepsilon)\to0$）；依分布 $F_n(x)\to F(x)$ 在连续点成立；几乎必然 $P(\lim X_n=X)=1$；$r$ 阶矩收敛 $E|X_n-X|^r\to0$。
\item 关系：$r$ 阶矩收敛 $\Rightarrow$ 依概率收敛 $\Rightarrow$ 依分布收敛；几乎必然收敛 $\Rightarrow$ 依概率收敛。反过来一般不成立，考试只考“谁推谁”和反例的名字。
\item 条件期望：离散情形 $E(X\mid Y=y)=\sum_x x\,P(X=x\mid Y=y)$；连续情形 $E(X\mid Y=y)=\int x f_{X\mid Y}(x\mid y)dx$，其中 $f_{X\mid Y}(x\mid y)=\frac{f(x,y)}{f_Y(y)}$。
\item 四条常用性质：线性；全期望公式 $E[E(X\mid Y)]=E(X)$；独立时 $E(X\mid Y)=E(X)$；$E[g(Y)X\mid Y]=g(Y)E(X\mid Y)$。
\item 条件方差公式：$\mathrm{Var}(X)=E[\mathrm{Var}(X\mid Y)]+\mathrm{Var}(E[X\mid Y])$，做“随机个随机变量之和”的方差时直接用它。
\end{itemize}
\textbf{完整例题（笔记题 5，也是期末常见题）}

设 $X_1,X_2,\dots$ 独立同分布，$N$ 是非负整值随机变量且与 $\{X_i\}$ 独立，$Y=\sum_{i=1}^{N}X_i$。求 $E(Y)$ 与 $\mathrm{Var}(Y)$。

第一步，对 $N$ 取条件：$E(Y)=E[E(Y\mid N)]$（全期望公式）。 第二步，给定 $N=n$ 时，$Y$ 是 $n$ 个独立同分布变量之和，$E(Y\mid N=n)=nE(X_1)$，所以 $E(Y\mid N)=NE(X_1)$。 第三步，再取期望：$E(Y)=E(N)E(X_1)$。这一步是 Wald 公式。 第四步，方差用条件方差公式：$\mathrm{Var}(Y)=E[\mathrm{Var}(Y\mid N)]+\mathrm{Var}(E[Y\mid N])$。其中 $\mathrm{Var}(Y\mid N)=N\,\mathrm{Var}(X_1)$，$E[Y\mid N]=N E(X_1)$。 第五步，代入得 $\mathrm{Var}(Y)=E(N)\mathrm{Var}(X_1)+\mathrm{Var}(N)[E(X_1)]^2$。 注意：第三步能提出 $E(X_1)$，靠的是 $N$ 与 $X_i$ 独立；若不独立，这条公式失效。

\Needspace{4\baselineskip}
\subsection{举一反三}
\PracticeQuestion{R4}{用上面的结果，设 $N\sim P(\lambda)$、$X_i\sim \mathrm{Exp}(\mu)$，求 $E(Y)$ 与 $\mathrm{Var}(Y)$ 的数值表达式。}{42}
\PracticeQuestion{V4}{若 $N$ 与 $X_i$ 不独立（例如 $N$ 由 $X_1$ 决定），指出 Wald 公式在哪一步失效，并举一个反例说明 $E(Y)\ne E(N)E(X_1)$。}{40}
\PracticeQuestion{M4}{判断：$E(X\mid Y)$ 是随机变量还是常数；当 $X$ 与 $Y$ 独立时它等于什么；并解释为什么 $E[E(X\mid Y)]=E(X)$ 不需要独立性。}{38}
\clearpage
\Needspace{6\baselineskip}
\section{五、独立练习}
先独立完成，再对答案。允许用 60 分钟。

\PracticeQuestion{A1}{设 $X$ 的分布函数为 $F(x)$，写出 $P(a<X\le b)$、$P(X=a)$ 的表达式，并说明为什么 $F$ 必须右连续。}{23}
\PracticeQuestion{A2}{设 $X\sim B(3,\frac12)$，求 $E(X)$、$\mathrm{Var}(X)$ 与 $P(X\ge 2)$。}{26}
\PracticeQuestion{A3}{求 $X\sim \mathrm{Exp}(\lambda)$ 的特征函数，并用它求 $E(X)$。}{34}
\PracticeQuestion{A4}{设 $X\sim N(\mu,\sigma^2)$，写出特征函数，并求 $E(X)$ 与 $\mathrm{Var}(X)$。}{35}
\PracticeQuestion{A5}{设 $X,Y$ 独立，$X\sim P(2)$、$Y\sim P(3)$，求 $X+Y$ 的分布与 $P(X+Y=1)$。}{35}
\PracticeQuestion{A6}{设 $X\sim B(n,p)$，用特征函数验证 $E(X)=np$。}{35}
\PracticeQuestion{A7}{设 $X_i$ 独立同分布、$N$ 与它们独立，$Y=\sum_{i=1}^N X_i$，写出 $E(Y)$ 与 $\mathrm{Var}(Y)$ 的公式（不代入具体分布）。}{35}
\PracticeQuestion{A8}{举一个“依概率收敛但不几乎必然收敛”的例子（可只描述构造思路）。}{35}
\clearpage
\Needspace{6\baselineskip}
\section{六、答案与解析}
\textbf{A1.} $P(a<X\le b)=F(b)-F(a)$；$P(X=a)=F(a)-F(a-0)$。右连续保证 $F(a)=F(a+0)$，从而 $P(X\le a)$ 与 $P(X<a)$ 的差恰好是单点质量，否则分布函数在跳点处会“丢掉”质量。

\textbf{A2.} $E(X)=1.5$，$\mathrm{Var}(X)=0.75$，$P(X\ge2)=\frac{4}{8}=\frac12$。

\textbf{A3.} $\varphi(t)=\left(1-\frac{it}{\lambda}\right)^{-1}$，$\varphi'(0)=\frac{i}{\lambda}$，所以 $E(X)=\frac{1}{\lambda}$。

\textbf{A4.} $\varphi(t)=e^{i\mu t-\frac12\sigma^2t^2}$，$E(X)=\mu$，$\mathrm{Var}(X)=\sigma^2$。

\textbf{A5.} $X+Y\sim P(5)$，$P(X+Y=1)=5e^{-5}$。

\textbf{A6.} $\varphi(t)=(pe^{it}+1-p)^n$，$\varphi'(t)=n(pe^{it}+1-p)^{n-1}ipe^{it}$，$\varphi'(0)=inp$，故 $E(X)=\frac{\varphi'(0)}{i}=np$。

\textbf{A7.} $E(Y)=E(N)E(X_1)$；$\mathrm{Var}(Y)=E(N)\mathrm{Var}(X_1)+\mathrm{Var}(N)[E(X_1)]^2$。

\textbf{A8.} 令 $X_n$ 独立、$P(X_n=1)=\frac1n$、$P(X_n=0)=1-\frac1n$，则 $X_n\xrightarrow{P}0$；但由 Borel–Cantelli，$X_n=1$ 无穷次发生，故 $X_n$ 不几乎必然收敛到 0。构造思路是“概率趋于 0 但可列次试验中几乎必然发生无穷多次”。

\Needspace{6\baselineskip}
\section{七、复习要点与反馈}
\begin{itemize}
\item 今天必须能不看书写出：分布函数三性质；矩母函数与特征函数的定义、导数与矩的关系；五个常见特征函数；条件期望四条性质；条件方差公式。
\item 做题时先写“设”和“所求”，再动笔算；特征函数题一定写“由唯一性定理”这一步。
\item 反馈方式：把 A1—A8 的答案拍照或打字发我，我会据此记录掌握状态、更新错题本，并在第 2 章课程里调整讲解密度。没有反馈前，本课所有条目保持“待验证”。
\end{itemize}
\end{document}
